\section{Colourings, obstructions, and counting}\label{co:color}

Throughout this section graphs are finite, simple and undirected. Write
$[n]=\{0,\ldots,n-1\}$. A proper $q$-colouring assigns elements of $[q]$ to
$[n]$ and gives the two ends of every edge different values. An independent set
has no internal edge. A clique has every possible internal edge. Let
$\chi,\alpha,\omega$ be the least number of colours, largest independent-set
size and largest clique size, respectively. All three extrema include witnesses.
Scan subset masks below $2^n$ for the last two. Scan $q\le n$ and the
$q^n$ words of length $n$ for the first, using $0^0=1$.
Giving each vertex its own colour certifies that this scan succeeds.

\begin{proposition}[Basic bounds and their witnesses]\leavevmode\label{co:greedy}
For every graph, $\omega\le\chi$ and $n\le\alpha\chi$.
If each vertex has at most $d$ earlier neighbours in a specified order,
a greedy algorithm produces a proper $(d+1)$-colouring. In particular
$\chi\le\Delta+1$, where $\Delta$ is the maximum degree, taken as zero on the
empty graph. If every nonempty induced subgraph has a vertex of degree at
most $d$, the algorithm also gives a $(d+1)$-colouring.
\end{proposition}
\begin{proof}
Colours on a clique are distinct. Each colour class is independent, so summing
its size gives the second inequality. At each step mark the colours of earlier
neighbours and pick the first unmarked colour among $0,\ldots,d$. At most
$d$ have been marked. The prefix of assigned vertices remains properly
coloured after the new vertex is inserted. For the last assertion repeatedly
remove the first vertex of current degree at most $d$ and colour in reverse
removal order. The input hypothesis is a scan over all $2^n$ induced
subgraphs, and the invariant for each of the $n$ removal steps checks the
residual vertex mask and its size.
\end{proof}

\begin{theorem}[Bipartition or odd cycle]\leavevmode\label{co:bip}
An algorithm returns either a proper $2$-colouring or a simple cycle of
odd length at most $n$. Exactly one of these outcomes is possible.
A bipartite graph with $c$ components has precisely $2^c$ proper
$2$-colourings.
\end{theorem}
\begin{proof}
Root a spanning tree in each component at its least vertex. Colour each
vertex by the parity of its tree distance from its root. Inspect every edge
$uv$. If the parities agree, follow the two parent chains to their last
common vertex $w$ on the paths from their root. The two tails after $w$ are
internally disjoint. The sum of their lengths is
$d(u)+d(v)-2d(w)$, which is even. The edge $uv$ closes an odd cycle.
This edge cannot be a tree edge, since adjacent levels have opposite parity.
The same construction works if $w=u$ or $w=v$. If no edge fails the test,
the colouring is proper. A proper $2$-colouring alternates along any walk,
so it cannot close an odd cycle. When the graph passes the test, the root
colour in each component forces every other colour along the unique tree
path. Either root choice works, since it flips both endpoints of every edge
in that component. These $c$ independent choices give the claimed count,
including the single empty colouring when $c=0$.
\end{proof}
Thus bipartite means $\chi\le 2$. A bipartite graph containing an edge
has $\chi=2$, a nonempty edgeless graph has $\chi=1$, and the empty graph
has $\chi=0$. Every induced subgraph of a bipartite graph is bipartite.
Its largest clique has size $2$, $1$ or $0$ in those same three cases,
so bipartite graphs are perfect: $\omega=\chi$ on every induced subgraph.
In particular $\chi(C_j)=2$ for even $j\ge 4$ and $\chi(C_j)=3$ for odd
$j\ge3$. For odd $j$, alternate on the first $j-1$ vertices and use a third
colour on the last.

\begin{proposition}[Cographs and optimal colourings]\leavevmode\label{co:cographs}
Given an expression built from single vertices by disjoint union and join,
one can produce an optimal colouring and a clique of that size on every
induced subgraph. Thus every cograph given by its construction is perfect.
\end{proposition}
\begin{proof}
A join adds all edges between its two disjoint operands. At a disjoint union,
a clique is contained in one side (the empty clique may be placed on either
side), and the sides may use the same palette. Take the larger of the two
cliques and reuse the palettes. At a join, concatenate the cliques and shift
the palette of the second side above that of the first: the size and palette
now both add. At a retained leaf use a one-element clique and one colour,
and at a deleted leaf use the empty clique and no colours. Evaluating the
expression from leaves to root therefore handles any given induced-subgraph
mask. On each node the two witnesses have matching objective values, so the
clique lower bound proves optimality there. A postorder list with children
pointing backwards permits one bounded induction through the expression.
\end{proof}

\begin{theorem}[Brooks]\leavevmode\label{co:brooks}
Every connected nonempty simple graph of maximum degree $\Delta$ that is
neither complete nor an odd cycle has a constructive $\Delta$-colouring.
\end{theorem}
\begin{proof}
In a connected graph of maximum degree at most $D$ with a vertex $z$ of
degree less than $D$, root a spanning tree at $z$ and greedily colour from
greatest depth down. A nonroot has an uncoloured parent, and the root has
fewer than $D$ neighbours. This uses at most $D$ colours. It settles the
nonregular case with $D=\Delta$. A connected regular graph of degree at
most two is a vertex, an edge, or a cycle: in degree two follow successive
edges without turning back until the first repeat, which closes a cycle
with no unused incident edge, hence spans by connectivity. The even cycle
is two-colourable. Assume henceforth that $G$ is regular of degree
$\Delta\ge3$ and is noncomplete.

If $w$ is a cut vertex, each component of $G-w$ together with $w$ has
$w$ of degree less than $\Delta$ there, since $w$ also has neighbours in
another component. Apply the rooted construction separately and permute
the palettes to give $w$ the same colour in all components.
Suppose now there is no cut vertex. We find nonadjacent neighbours $x,y$
of some $z$ such that $G-\{x,y\}$ is connected. If every pair can be
deleted while preserving connectivity, take the first and third vertices
of a shortest path between any nonadjacent pair, with $z$ in
the middle. Otherwise choose $z,w$ whose joint deletion disconnects.
Write $J=G-z$, which is connected. In each component $C$ of $J-w$ there
is a neighbour of $z$, since $G-w$ is connected. Choose such a neighbour
$x_C$ of greatest distance from $w$ measured in $J$. Then $J-x_C$ is
connected. Indeed, a component $D$ of $J-x_C$ avoiding $w$ would lie
in $C$, since every other component of $J-w$ can reach $w$ without
$x_C$. Connectivity of $G-x_C$ forces an edge from $z$ to some $d\in D$.
Every $w$--$d$ path in $J$ uses $x_C$, making $d$ strictly farther from
$w$ than $x_C$, a contradiction. Pick $x,y$ of this kind in two
different components. They are nonadjacent. In each affected component all
remaining vertices reach $w$ avoiding its deleted vertex, and the initial
path to $w$ stays inside that component. Thus $J-\{x,y\}$ is connected.
Vertex $z$ still has a third neighbour, proving the assertion for $G$.

Give $x,y$ the same colour. In $G-\{x,y\}$ root a spanning tree at
$z$ and colour in decreasing depth. Every vertex other than $z$ has an
uncoloured parent, and at $z$ the already coloured neighbours $x,y$
forbid just one colour between them. A choice among $\Delta$ colours
is available at every step. Vertex pairs, components and shortest
path lengths in this proof are all found by scans on the fixed input.
\end{proof}

\subsection{The chromatic polynomial and acyclic orientations}
For $F\subseteq E$ let $c(F)$ be the number of components of the spanning
graph $([n],F)$, counting isolated vertices. Define the signed coefficient
list of a polynomial by a finite sum
\[
 P_G(X)=\sum_{F\subseteq E}(-1)^{|F|}X^{c(F)}.
\]
Each coefficient is a difference of two natural counts at most $2^{|E|}$.
Polynomial statements below mean finite coefficient-list identities.

\begin{theorem}[Chromatic identities]\leavevmode\label{co:poly}
$P_G(q)$ counts proper $q$-colourings for each natural $q$. It is monic
of degree $n$ and, for $n\ge1$, its coefficient of $X^{n-1}$ is $-|E|$.
For an edge $e=uv$, deleting $e$ gives $D$ and identifying its ends in $D$
and suppressing duplicate edges gives $H$. Then
$P_G=P_D-P_H$. Disjoint unions multiply chromatic polynomials.
\end{theorem}
\begin{proof}
For a fixed colouring word $a$, let $B(a)$ consist of its monochromatic
edges. Its proper-colouring indicator is
$\sum_{F\subseteq B(a)}(-1)^{|F|}$: when $B(a)$ is nonempty, toggling its
least edge pairs terms of opposite sign, and otherwise just the empty
subset occurs. Sum over words and interchange the two finite sums. A word
constant on the edges of $F$ is specified by one value per component,
giving $q^{c(F)}$ possibilities. The only spanning subgraph with $n$
components has no edge. A spanning simple subgraph with $n-1$ components
has exactly one doubleton component, carrying one edge, and all other
vertices isolated. This gives the first two coefficients.

For deletion and contraction, the proper words on $D$ split by whether the
colours at $u,v$ agree. The unequal words are precisely the proper words on
$G$. An equal word descends uniquely to $H$, since the common colour is the
one on the identified vertex, and duplication of an inequality has no
effect. The two words on the sides of a disjoint union can be chosen
independently, so the product identity also holds on natural inputs.
To pass from evaluations to coefficients, a nonzero integer polynomial of
degree $h$ cannot vanish on $h+1$ distinct integers: at a root $b$,
$P(X)-P(b)$ is $(X-b)$ times the polynomial obtained from
$X^j-b^j=(X-b)\sum_{i<j}X^{j-1-i}b^i$. At every other root that quotient
also vanishes by integer cancellation. Iterating gives a nonzero constant
with a root, an impossibility. Only evaluations at $0,\ldots,n$ are
needed here. For an explicit bound in this root argument, if the initial
absolute coefficients are at most $M$ and the roots have magnitude at most
$N$, each division raises the coefficient bound by a factor at most
$(h+1)(N+1)^h$. Iterate this factor at most $h$ times.
\end{proof}
For the complete graph $K_n$ this gives $P_{K_n}(q)=(q)_n$, the number of
injections from $[n]$ to $[q]$. It is zero for $q<n$. For a forest with
$c$ components, orient a rooted tree in each component away from its root.
For $q\ge 1$ the roots have $q$ choices each and each other vertex has
$q-1$ choices after its parent. Thus $P_G(X)=X^c(X-1)^{n-c}$ by the
same coefficient argument, evaluating at $1,\ldots,n+1$.

\begin{theorem}[Stanley's evaluation]\leavevmode\label{co:stanley}
The number $a(G)$ of acyclic orientations of $G$ is $(-1)^nP_G(-1)$.
\end{theorem}
\begin{proof}
Keep $D,H,u,v$ as above. In an acyclic orientation of $D$ there is at
most one direction of reachability between $u$ and $v$. If a directed
path joins them, orienting $e$ in the same direction is the only acyclic
extension. If neither is reachable from the other, both orientations extend.
The latter orientations of $D$ are in bijection with acyclic orientations
of $H$. Indeed, at a common neighbour the two edges incident to $u,v$ have
the same direction relative to that neighbour, since opposite directions
would make a two-step path. Thus the orientation descends. A directed
simple cycle through the identified vertex would lift to a directed cycle
or a directed path between $u,v$. A cycle avoiding that vertex would already
lie in $D$. Conversely the lift of an acyclic orientation of $H$ allows
neither a directed cycle nor a path between $u,v$, because either would
map to a nonempty directed closed walk. A shortest nonempty directed
closed subwalk is a directed cycle. The edges of $D$ all map to actual
edges of $H$, so the image walk has positive length.
Consequently $a(G)=a(D)+a(H)$. The signed evaluation has the same
recurrence because $H$ has one fewer vertex, and both quantities are one
on edgeless graphs. Induct on the edge count bounded by the initial edge
count, bounding all smaller graph codes by the fixed initial order.
\end{proof}

\subsection{Chordal graphs and perfect elimination}
A graph is chordal when every simple cycle of length at least four has a
chord, an edge joining two nonconsecutive vertices. A vertex is simplicial
when its neighbours form a clique.

\begin{theorem}[Minimal separators and elimination]\leavevmode\label{co:chordal}
In a chordal graph every inclusion-minimal vertex separator for a
nonadjacent pair is a clique. Every noncomplete chordal graph has two
nonadjacent simplicial vertices. A graph is chordal precisely when it has
an ordering in which the later neighbours of every vertex form a clique.
Every chordal graph is perfect.
\end{theorem}
\begin{proof}
For nonadjacent $x,y$, choose an inclusion-minimal subset
$S\subseteq V\setminus\{x,y\}$ separating them. Such a set exists by
deleting all other vertices and scanning subsets. Let $A,B$ be the
components of $x,y$ after deleting $S$. Every $s\in S$ has a neighbour
in each component: a simple $x$--$y$ path on restoring $s$ must pass
through $s$, with the portions on either side contained in $A,B$. If
$p,q\in S$ were nonadjacent, there would be a $p$--$q$ path through $A$
and one through $B$. Take shortest ones with interiors in those respective
components. Neither path has a chord by minimality, and there is no edge
between the two interiors. Their union is an induced cycle of length at
least four, a contradiction. This also covers $S=\varnothing$.

Use induction on the size of an induced subgraph of the fixed input.
For a noncomplete graph choose $x,y,S,A,B$ as above. In the smaller
chordal induced graph on $A\cup S$, find a simplicial vertex in $A$:
if this graph is complete, choose any vertex of $A$. If it is noncomplete,
the induction hypothesis produces two nonadjacent simplicial vertices,
at most one of them in the clique $S$. All neighbours of the chosen
vertex in the whole graph already belong to $A\cup S$, so it remains
simplicial. Repeating with $B$ gives the second vertex, nonadjacent to
the first. The complete case has every vertex simplicial.
Repeated removal therefore produces the required ordering. Conversely, on
an induced cycle of length at least four, the first vertex in such an
ordering would have two nonadjacent later neighbours. Finally colour in
reverse order. Each vertex together with its later neighbours is a clique,
so at most $\omega-1$ colours are forbidden. Every induced subgraph is
chordal by inspection of its cycles, hence $\chi=\omega$ throughout.
\end{proof}

\subsection{Colouring edges by fans}
In a proper edge colouring, edges sharing an endpoint have distinct colours.
Equivalently this is a vertex colouring of the line graph, whose vertices
are the original edges and whose adjacencies mean a shared endpoint.

\begin{theorem}[Vizing]\leavevmode\label{co:vizing}
Every simple graph of maximum degree $\Delta$ has an edge colouring with
$\Delta+1$ colours. Its least possible palette size is $\Delta$ or
$\Delta+1$.
\end{theorem}
\begin{proof}
Use a fixed palette of $\Delta+1$ colours. Extend a proper colouring of
the first $t$ edges to include the next edge $xy$. Ignore later edges.
A colour missing at a vertex occurs on none of its coloured incident
edges. At least one colour is missing at every vertex.
Make a list of distinct neighbours $w_0=y,w_1,\ldots,w_k$ of $x$.
Append a neighbour $w_j$ whenever the coloured edge $xw_j$ has a colour
missing at an earlier $w_{p(j)}$, recording one such predecessor $p(j)<j$.
Stop when no new neighbour can be appended.

For any target index $i$, follow predecessors down to $0$, then reverse
the chain. Shifting the colour of each successive spoke $xw_j$ to the
preceding spoke $xw_{p(j)}$ leaves the target spoke empty. Each moved
colour was missing at its new outer endpoint. At $x$ the distinct
colours are permuted along with the single empty slot. We call this
operation rotation. It succeeds whenever the recorded missing colours on
that chain still are missing.

Choose $\alpha$ missing at $x$ and $\beta$ missing at $w_k$.
If $\beta$ is missing at $x$, rotate to $k$ and use $\beta$ there.
Otherwise the $\beta$-spoke is $xw_j$ for some $1\le j\le k$, by
maximality of the list, and $\alpha\ne\beta$. Set $i=p(j)$. Both $w_i$
and $w_k$ miss $\beta$ and they are distinct. In the graph of
$\alpha$- and $\beta$-edges every vertex has degree at most two. The
component of $x$, which has degree one there, is a path. Iterating from its
only edge cannot revisit a vertex (a first repeat would require a third
incident edge or a second edge at $x$). Thus at most one of $w_i,w_k$
belongs to that component, since both have degree at most one. Swapping
$\alpha$ and $\beta$ throughout any component preserves propriety.

If $w_i$ is outside the component of $x$, swap on its component. This
makes $\alpha$ missing at $w_i$. Rotation to $i$ uses only spokes
of index at most $i<j$. None has colour $\alpha$ or $\beta$.
All its recorded missing colours survive the swap. Rotate and use
$\alpha$ at $xw_i$. In the remaining case swap on the component of
$w_k$, which contains neither $w_i$ nor $x$. This makes $\alpha$ missing
at $w_k$. On the chain to $k$, the only possible spoke of colour $\beta$
is $xw_j$, and its recorded predecessor $w_i$ still misses $\beta$.
There are no $\alpha$-spokes at $x$. Again the rotation is valid and
$\alpha$ fills the final empty spoke. Each extension traverses at most
$n$ fan vertices and one two-colour component. The outer induction checks
propriety of the coloured edge prefix. Finally the incident edges of a
maximum-degree vertex already require $\Delta$ colours. The empty edge
set has least palette size zero.
\end{proof}

A further vertex-colouring certificate is an acyclic orientation whose
longest directed path has at most $q$ \emph{vertices}. Direct each edge
from smaller to larger colour. Colours strictly increase along a directed
path. In the reverse direction, colour $v$ by the maximum number of
vertices on a directed path ending at $v$, minus one. A longest such path
ending at $u$ extends across an arc $u\to v$, since otherwise $v$ would
already occur on it and give a directed cycle. The colour therefore
strictly increases on each arc. Paths are found by scanning words of at
most $n$ vertices. The least such longest-path order, over acyclic
orientations, is exactly $\chi$ (the longest-path order of the empty graph
is zero).

\section{Extremal graph theory with explicit witnesses}\label{co:extremal}
\subsection{The exact clique-free maximum}
Write $T_r(n)$ for the complete $r$-partite graph on $n$ vertices whose
part sizes differ by at most one, allowing empty parts. Thus, for $r\ge1$
and $n=ar+b$ with $0\le b<r$, its part sizes are $b$ copies of $a+1$ and
$r-b$ copies of $a$. Let $t_r(n)$ denote its edge count.
\begin{theorem}[Tur\'an, including equality]\leavevmode\label{co:turan}
A graph on $n$ vertices containing no $(r+1)$-clique has at most
\[
 t_r(n)=\binom n2-b\binom{a+1}2-(r-b)\binom a2
\]
edges. Equality holds precisely for graphs isomorphic to $T_r(n)$. A
partition giving this isomorphism can be constructed. Strictly more than
$t_r(n)$ edges yields an explicit $(r+1)$-clique.
\end{theorem}
\begin{proof}
First construct a complete $r$-partite graph on the same vertex set having
at least as many edges, with equality only if the two graphs actually have
the same edges. For $r=1$ the input is edgeless. For $r>1$, an empty graph
of order zero needs only empty parts. Otherwise take a vertex of maximum
degree $d$, put $S$ equal to its neighbours and $T=V\setminus S$.
The graph on $S$ contains no $r$-clique, since that clique together with the
chosen vertex would be forbidden. Let $e(S),e(T),e(S,T)$ be the internal and
crossing edge counts. Summing degrees over $T$ gives
\[
 e(S,T)+2e(T)\le d|T|,\qquad
 e(G)\le e(S)+d|T|-e(T)\le e(S)+|S||T|.
\]
Recursively replace $G[S]$ by a complete $(r-1)$-partite graph, and
make $T$ an independent part joined to all of $S$. These inequalities prove
the promised domination. Equality forces $e(T)=0$, every one of the
$|S||T|$ cross edges present, and equality inside $S$. Hence the
replacement made no change. This only requires $r$ successive passes.
Alternatively, for induction on $r$ fix a bound $N$ on vertex labels
and quantify over masks of $[N]$ and matrices below $2^{N^2}$ inside
the induction formula.

A complete multipartite graph has $\binom n2-\sum_i\binom{a_i}2$ edges.
If two parts have sizes $x\ge y+2$, moving a vertex from the larger to the
smaller increases the edge count by $(x-1)-y\ge1$. There can be at most
$\binom n2$ improving moves. At termination all sizes are $a$ or $a+1$,
with the stipulated multiplicities. This proves the formula and the equality
condition. The extremal construction forbids an $(r+1)$-clique by the
pigeonhole principle. For a graph above the bound, scan all subset masks
for a clique of the required size. Exhaustion contradicts the proved bound.
\end{proof}
For $r=2$ this is Mantel's theorem, with bound
$\lfloor n/2\rfloor\lceil n/2\rceil=\lfloor n^2/4\rfloor$ and the unique
balanced complete bipartite extremizer up to isomorphism.
For later use,
$t_r(n)\le (1-1/r)n^2/2$: the identity
$r\sum_i a_i^2-(\sum_i a_i)^2=\sum_{i<j}(a_i-a_j)^2$
gives the bound from $2t_r(n)=n^2-\sum_i a_i^2$.

\subsection{Dense graphs contain a prescribed multipartite graph}
The following bound gives the Erd\H{o}s--Stone theorem its finite,
witness-producing formulation. All densities in this and the next
subsection abbreviate inequalities of naturals after clearing positive
denominators.

\begin{lemma}[A dense relation contains a box]\leavevmode\label{co:box}
Let $k,s,d\ge1$ and let $R\subseteq[n]^k$ satisfy $d|R|\ge n^k$.
Define $d_0=d$ and
$d_{j+1}=2d_j(4d_j)^s$. If $n\ge 4s d_k$, there are $s$-element
sets $A_1,\ldots,A_k$ with $A_1\times\cdots\times A_k\subseteq R$.
\end{lemma}
\begin{proof}
With current dimension $l\ge2$ and density denominator $D$, each of the
$n^{l-1}$ prefixes has a number of extensions in $[n]$. At least
$n^{l-1}/(2D)$ prefixes have at least $n/(2D)$ extensions: the remaining
prefixes together account for at most $n^l/(2D)$ tuples and each retained
prefix accounts for at most $n$. From each retained prefix one can choose
an ordered list of $s$ distinct extensions in at least $(n/(4D))^s$ ways,
because fewer than $s$ previously used labels are excluded and $n\ge4Ds$.
Sum over prefixes and reorder the count by extension list. There are at
most $n^s$ lists. Thus some list has at least
$n^{l-1}/[2D(4D)^s]$ common prefixes. It is found by a bounded scan.
Its entries form $A_l$, and the relation of common prefixes is the input
for dimension $l-1$. At dimension one there are at least $n/D\ge s$
members, so take the first $s$. The invariant after each projection says
that every current tuple concatenated with any choices from the recorded
suffix lies in $R$, and that the current relation meets its cardinality
bound with the new $D$. All relation masks have at most $n^k$ bits.
\end{proof}

\begin{theorem}[Effective Erd\H{o}s--Stone]\leavevmode\label{co:es}
Fix $r\ge2$, $s\ge1$ and a positive integer $h$. Write
$\lambda=1-1/(r-1)$, $q=\max\{r,2h+1\}$, and
$d=2^{r+1}h\binom qr$. Use Lemma~\ref{co:box} with this $d$ and set
\[
 N(r,s,h)=\max\{q,2r,4s d_r\}.
\]
If $n\ge N(r,s,h)$ and $e(G)\ge(\lambda+1/h)\binom n2$,
one can find $r$ disjoint sets of size $s$ with every edge between
different sets present.
\end{theorem}
\begin{proof}
The average edge count of an induced $q$-vertex subgraph equals
$e(G)\binom q2/\binom n2$, since each edge lies in
$\binom{n-2}{q-2}$ such subsets. If a $q$-set has no $r$-clique,
Tur\'an's bound gives at most
\[
 \lambda q^2/2=
 \bigl(\lambda+\lambda/(q-1)\bigr)\binom q2
 \le (\lambda+1/(2h))\binom q2
\]
edges. All other sets have at most $\binom q2$. Subtracting the bound
for the former from the average shows that at least a $1/(2h)$
fraction of all $q$-sets contain an $r$-clique: if their fraction is $f$,
the average divided by $\binom q2$ is at most
$\lambda+1/(2h)+f$. Every $r$-clique lies in
$\binom{n-r}{q-r}$ such sets. Counting incidences shows that the number
of $r$-cliques is at least $\binom nr/(2h\binom qr)$. Here one uses
$\binom nq\binom qr=\binom nr\binom{n-r}{q-r}$.

Let $R$ be the relation of ordered $r$-cliques. Its size is $r!$ times
that count, and $(n)_r\ge (n/2)^r$ for $n\ge 2r$, so $d|R|\ge n^r$.
Apply Lemma~\ref{co:box}. Coordinates in a member of $R$ are distinct,
so the $r$ resulting sets must be pairwise disjoint: an element shared by
two sets could be chosen in those two coordinates simultaneously. Choosing
any pair from two different sets and completing the remaining coordinates
proves that all required edges are present.
\end{proof}

In particular, given a graph $H$ with $\chi(H)=r\ge2$, take $s=|V(H)|$.
Its optimal colouring gives an explicit embedding into this complete
multipartite subgraph, by assigning its same-colour vertices distinct
labels in the corresponding part. In the other direction, every
$(r-1)$-partite graph avoids $H$, since an embedding would colour $H$ with
$r-1$ colours. Such a graph on $n\ge2$ can attain density at least
$\lambda$: among all $(r-1)^n$ assignments of vertices to labelled parts,
the fraction separating any fixed edge of $K_n$ is $\lambda$, so one
assignment cuts at least $\lambda\binom n2$ pairs. Fill all those cross
edges. Thus for $n\ge N(r,s,h)$ the maximum edge density of an
$H$-free graph lies in $[\lambda,\lambda+1/h)$. This is a
primitive-recursive modulus for the usual density limit.

\subsection{The finite graph Ramsey theorem}
\begin{theorem}[Two colours and the binomial bound]\leavevmode\label{co:ramsey}
For $a,b\ge1$ put $B(a,b)=\binom{a+b-2}{a-1}$. Every graph on at
least $B(a,b)$ vertices has an $a$-clique or a $b$-vertex independent
set. A bounded algorithm returns one. In particular the two-colour Ramsey
number $R(a,b)$ exists and satisfies $R(a,b)\le B(a,b)$.
\end{theorem}
\begin{proof}
With $a=1$ or $b=1$, use any vertex. Otherwise fix a vertex $v$. Among the
remaining vertices either there are at least $B(a-1,b)$ neighbours, or
there are at least $B(a,b-1)$ non-neighbours. If both tests failed, the two
sets would together have at most $B(a,b)-2$ elements by Pascal's identity.
In the first branch restrict to the neighbours and lower $a$ by one.
An independent-set answer works as it stands, and a clique answer extends
by $v$. In the second branch restrict to the non-neighbours and lower $b$
by one. Now extend just the independent-set answer by $v$. A stored stack
of at most $a+b$ branch choices implements the descent and reconstruction.
One may induct on the sum of the targets with all varying targets bounded
by the initial $a,b$ and all varying vertex sets masks on the original
ground set. The least working graph order is found by checking all simple
graph codes at each order through $B(a,b)$. It works at all larger orders
by restriction to that many vertices.
\end{proof}
For example $B(3,3)=6$. The $5$-cycle has no triangle and its complement
is a $5$-cycle (traverse the vertices in steps of two), so it has no
independent triple either. Hence $R(3,3)=6$. For $k\ge1$ edge colours
and target $s\ge1$, set $L_1=s$ and $L_{i+1}=B(s,L_i)$. On $L_{i+1}$
vertices group the edges of one colour against all other edges. The
two-colour result returns an $s$-set of the first colour or an $L_i$-set
on which the remaining $i$ colours suffice. Iterating yields a
monochromatic $s$-set. The nested values $L_i$ are given by ordinary
primitive recursion on $i$.

\subsection{An equitable regularity lemma with an arithmetic bound}
The following proof supplies the exceptional set, all equal-size classes,
and the witnesses used to refine them. For disjoint nonempty sets $A,B$
write $d(A,B)=e(A,B)/(|A||B|)$. A pair is $\varepsilon$-regular if
$|d(X,Y)-d(A,B)|\le\varepsilon$ whenever $X\subseteq A$, $Y\subseteq B$,
$|X|\ge\varepsilon|A|$, and $|Y|\ge\varepsilon|B|$.

\begin{theorem}[Finite Szemer\'edi regularity]\leavevmode\label{co:regularity}
Let $h\ge2$, $k_0\ge1$, and $\varepsilon=1/h$.
Put $J=4h^5+1$, $K_0=\max\{k_0,2hJ\}$,
$K_{i+1}=4K_i4^{K_i}$ and
$N=\max\{2hK_J,8K_J4^{K_J}\}$.
Every graph of order $n\ge N$ has a partition into an exceptional set
$V_0$ of size at most $n/h$ and $k$ nonempty equal-size classes, with
$k_0\le k\le K_J$. At most $k^2/h$ ordered pairs of distinct classes
are irregular. The partition and a bounded regularity test are
primitive recursive.
\end{theorem}
\begin{proof}
Use ordered pairs of vertices throughout, counting a diagonal entry as
zero. An auxiliary partition $\mathcal P$ consists of the main classes
and the exceptional vertices as individual singleton cells. Define its
energy by
\[
 Q(\mathcal P)=\sum_{A,B\in\mathcal P}\frac{|A||B|}{n^2}
                  d(A,B)^2.
\]
Here the density for $A=B$ also uses ordered pairs and denominator
$|A|^2$. We have $0\le Q\le1$. In a refinement, an old rectangle with
density $d$ is split into rectangles with relative areas $w_i$ and densities
$d_i$. Since $\sum_iw_i=1$ and $\sum_iw_i d_i=d$, its gain in mean
square is $\sum_iw_i(d_i-d)^2$. In particular the energy never decreases.

Suppose the main classes have size $m$ and a pair $A,B$ is irregular.
Scan the two subset masks to choose witnesses $X,Y$. If a refinement
respects those masks, $X\times Y$ is a union of its rectangles. Their
contribution to the mean square gain on $A\times B$ is at least
\[
 \frac{|X||Y|}{m^2}\bigl(d(X,Y)-d(A,B)\bigr)^2> h^{-4}.
\]
Indeed the weighted average of $d_i-d(A,B)$ on that subrectangle is
$d(X,Y)-d(A,B)$, and the average of squares dominates the square
of the average by the same variance identity. Use a single witness
pair per unordered irregular pair, transposing for the opposite order.
Inside each main class split by membership in all its chosen witness sets.
There are at most $2^k$ nonempty atoms per class. If more than $k^2/h$
ordered pairs are irregular and $km\ge n/2$, the increase in energy
is greater than $(km/n)^2 h^{-5}\ge1/(4h^5)$.

It remains to ensure equal sizes without undoing this gain. Start with
$K_0$ consecutive blocks of $\lfloor n/K_0\rfloor$ vertices and place
the remainder in $V_0$. At a refinement, divide every atom in vertex order
into blocks of a common size $m'=\lfloor m/4^k\rfloor$, transferring the
unused vertices of that atom to $V_0$. As individual singletons these
vertices also refine the atom. At most $k2^k m'\le n/2^k$ vertices are
added to $V_0$. Initially $|V_0|<K_0\le n/(2h)$, and $J$ such
additions total at most $Jn/2^{K_0}\le n/(2h)$ provided $k\ge K_0$.
These estimates propagate together: at stage $i<J$, suppose
$K_0\le k\le K_i$ and $|V_0|\le n/(2h)+in/2^{K_0}\le n/h$.
Then $km\ge n/2$, and the choice of $N$ gives $m\ge2\cdot 4^k$.
Consequently $m'\ge m/(2\cdot 4^k)\ge1$. Every old class contributes
at least one new main block since fewer than $2^k m'\le m/2^k<m$
of its vertices are discarded. The new class count is at most
$n/m'\le 4k4^k\le K_{i+1}$. This proves the required induction estimates.

The scan stops when the irregularity count satisfies the theorem.
Otherwise $J$ failed tests and refinements would increase energy by more
than $J/(4h^5)>1$. This proves termination with the stated bound.
An energy may be stored over the common denominator $n^2(n!)^2$ because
each ordered cell product $|A||B|$ divides $(n!)^2$. Cross-multiplication
therefore makes the energy-increment invariant a Boolean check, along with
the partition tests. The recurrence for the number of cells is uniform in
both parameters.
\end{proof}

\subsection{Large girth and large chromatic number}
\begin{theorem}[Erd\H{o}s, with a finite search bound]\leavevmode\label{co:girth}
For $g\ge3$ and $k\ge1$ there is a graph with no cycle of length at most
$g$ and with $\chi>k$. Its order is at most
$[64(g+1)k^2]^{2g}$, and such a graph can be produced by primitive
recursion in $g,k$.
\end{theorem}
\begin{proof}
Set $q=64(g+1)k^2$, $n=q^{2g}$, $D=n/q$, and $t=n/(2k)$.
Consider all $D^{\binom n2}$ words assigning one digit below $D$ to each
unordered pair of vertices, making it an edge exactly when its digit is
zero. An average always means the sum over this finite list divided by
its length. Prescribing $j$ particular edges to be present selects a
fraction $D^{-j}$ of the words by counting free digits.

Let $C$ count simple cycles of lengths $3$ through $g$.
Represent each cycle by its lexicographically first traversal among its
rotations and reversals. There are at most $n^j$ candidate traversals of
length $j$. Their $j$ distinct edges occur in a fraction $D^{-j}$ of words,
so the average of $C$ is at most $\sum_{j=3}^g q^j\le gq^g\le n/8$.
At most a quarter of the words have $C\ge n/2$, since summing $C$ on
those words alone gives at least their number times $n/2$.

Let $I$ count independent $t$-sets. Each candidate needs its
$\binom t2$ pairs absent, so the average of $I$ is at most
$2^n(1-1/D)^{\binom t2}$. Since $D\ge2$,
$(1+1/(D-1))^D\ge1+D/(D-1)\ge2$ by the binomial expansion.
Thus $(1-1/D)^D\le1/2$. Here $t\ge2$ and
\[
 \binom t2\ge t^2/4=n^2/(16k^2)\ge D(n+2),
\]
using $q\ge32k^2$ and $n+2\le2n$. The average of $I$ is at most
$1/4$, so at most a quarter of the words have $I\ge1$.
Consequently a word with $I=0$ and $C<n/2$ is found by bounded scan.
Choose one vertex of each of its short cycles, form the union of these
choices, and delete it. At least $n/2$ vertices remain and every short
cycle was hit. Every independent set remaining has size $<t$.
With only $k$ colours the graph could have at most $k(t-1)<n/2$
vertices, a contradiction.
\end{proof}

\section{Hamiltonian routes in graphs and tournaments}\label{co:hamilton}
A Hamiltonian cycle is a simple cycle visiting every vertex. The graph
has at least three vertices in the cycle results below.

\begin{theorem}[Closure, Ore, and Dirac]\leavevmode\label{co:ore}
If $u,v$ are nonadjacent in a simple graph of order $n\ge3$ and
$\deg(u)+\deg(v)\ge n$, adding $uv$ preserves and reflects existence
of a Hamiltonian cycle. Consequently a Hamiltonian cycle can be
constructed if every nonedge satisfies this inequality (Ore), in
particular if $2\deg(v)\ge n$ at every vertex (Dirac).
\end{theorem}
\begin{proof}
An existing cycle survives edge addition. Given a cycle after addition
that uses the new edge, deleting it gives a spanning path
$x_1=u,x_2,\ldots,x_n=v$ in the original graph. In the $n-1$ slots
$i=1,\ldots,n-1$ mark one type when $v$ is adjacent to $x_i$ and
a second type when $u$ is adjacent to $x_{i+1}$. The counts sum to
at least $n$, so some slot has both marks. Since $uv$ was a nonedge
this slot has $1<i<n-1$. Follow the path from $x_1$ to $x_i$, cross
to $x_n$, follow the path backwards to $x_{i+1}$, and close to $x_1$.
This gives the required cycle in the original graph. For Ore's condition,
add nonedges in order until the complete graph is obtained. Degree sums
only increase. Its vertex order gives a Hamiltonian cycle. Undo the
additions in reverse order by the preceding operation. There are at most
$\binom n2$ steps and each stored cycle is a permutation of $[n]$.
Dirac's condition implies Ore's by addition.
\end{proof}

\begin{theorem}[R\'edei and Camion]\leavevmode\label{co:tournament}
A tournament (exactly one arc on each pair of distinct vertices) has
a directed Hamiltonian path. A tournament on $n\ge3$ has a directed
Hamiltonian cycle if and only if it is strongly connected.
\end{theorem}
\begin{proof}
Insert the next vertex $x$ into an existing path immediately before the
first vertex to which $x$ has an outgoing arc, appending $x$ if there is
none. The predecessor, if any, has an arc to $x$. Starting with the
empty list gives a directed Hamiltonian path (the empty list is the
output at order zero).

A directed Hamiltonian cycle gives directed reachability in both directions
between each pair. For the converse, strong connectivity gives a directed
cycle: take an arc and a shortest path returning from its head to its tail.
Choose a longest directed simple cycle $C$ by scanning lists of length at
most $n$. If an outside vertex $x$ has incidences both ways with $C$,
a cyclic change of type on $C$ gives successive vertices $c,d$ with
$c\to x\to d$. Insert $x$ there to extend $C$. Thus any outside vertex
is either dominated by every vertex of $C$ (type $U$) or dominates every
vertex of $C$ (type $W$). If there is an outside vertex there must be
both types: otherwise either no arc leaves $C$ or no arc enters it.
There must be an arc $u\to w$ from $U$ to $W$, since the only possible
exit from $U$ is into $W$ and vertices of $U$ must reach $C$. Replace
any arc $c\to d$ of $C$ by $c\to u\to w\to d$, contradicting maximality.
The bounded scan therefore produces a spanning cycle.
\end{proof}
